\section{Introduction}
\label{sec:intro}
The parton distribution functions (PDFs) encode important nonperturbative information of strong interactions. Based on QCD factorization \cite{Collins:1989gx}, PDFs have been successfully extracted from high-energy scattering data with a good precision \cite{Lin:2017snn}. However, from both theoretical and practical points of view, extracting PDFs from first principle lattice QCD (LQCD) calculations must be done for testing non-perturbative sector of QCD, as well as needed for studying partonic structure of hadrons that could be difficult to do scattering experiments with.

Calculating PDFs in Euclidean-space LQCD {\it directly}, if not impossible, is difficult due to the time-dependence of the operators defining PDFs \cite{Lin:2017snn}.  A novel approach was suggested by Ji \cite{Ji:2013dva}, who introduced a set of LQCD-calculable quasi-PDFs and argued that the quasi-PDFs of hadron momentum $P_z$ become corresponding PDFs when $P_z$ is boosted to infinity. A number of other approaches to extract PDFs from LQCD calculations were also proposed~\cite{Liu:1993cv,Liu:1999ak,Liu:2016djw,Orginos:2017kos,Chambers:2017dov}. In Refs.~\cite{Ma:2014jla,Ma:2017pxb}, two of us proposed a QCD factorization based general approach to calculate PDFs in LQCD {\it indirectly}.
Similar to the extraction of PDFs from experimental data of factorizable and measurable hadronic cross sections, we proposed to
extract PDFs by the global analysis of data generated by LQCD calculations of good ``lattice cross sections" (LCSs), which are defined as hadronic matrix elements that satisfy (1) calculable in Euclidean-space LQCD, (2) renormalizable for ultraviolet (UV) divergences to ensure a reliable continue limit, and (3) factorizable to PDFs with infrared-safe matching coefficients. It is the (3)-factorization that relates the desired PDFs to the LQCD calculable LCSs.

To extract the rich, precise and flavor separated information on PDFs, it is necessary to find as many good LCSs as possible, since different flavor PDFs are likely to contribute to the same LQCD calculated LCSs.  For constructing good LCSs, we studied two types of operators in terms of (a) correlation of two gauge-dependent field operators with proper gauge links, which we call quasi-parton operators  since they cover all operators defining quasi-PDFs and more, and (b) correlation of two gauge-invariant currents. In our approach, LQCD calculation of each of these good LCSs in coordinate space provides the needed information to {\it constrain} the PDFs, similar to the role of measuring various experimental cross sections to {\it constrain} the PDFs. For a comparison, in Ji's proposal \cite{Ji:2013dva}, one focuses on reproducing PDFs by corresponding quasi-PDFs calculated in LQCD at a large enough $P_z$.



For LQCD calculations, it is relatively less expensive to calculate the type (a) operators comparing with type (b) operators, while the renormalization of the type (b) operators is much more simpler than that of the type (a) operators.
Renormalization of the type (b) operators is almost trivial, for which one only needs to renormalize the gauge-invariant local currents, which is well-known.  On the other hand, the renormalization of the type (a) operators is nontrivial due to the nonlocality of corresponding operators and potential power divergences.  QCD factorization of both types of operators have been studied in Ref.~\cite{Ma:2017pxb}, in which we found that multiplicative renormalizability of these operators is a necessary condition for the collinear factorization to be valid.

A lot of efforts have been devoted to explore the UV structure of the type (a) quasi-quark operators \cite{Ishikawa:2016znu,Chen:2016fxx,Monahan:2016bvm,Briceno:2017cpo, Xiong:2017jtn, Constantinou:2017sej, Alexandrou:2017huk, Chen:2017mzz}.
All-order multiplicative renormalizability of quasi-quark operators has been proved using two different methods: one relies on the auxiliary field technique \cite{Ji:2017oey,Green:2017xeu}, and the other is based on diagrammatic expansion \cite{Ishikawa:2017faj}. These proofs provided a firm theoretical basis for extracting the combination of quark distributions that are not sensitive to gluons from LQCD calculated  hadronic matrix elements of quasi-quark operators \cite{Lin:2014zya,Alexandrou:2015rja,Chen:2016utp,Alexandrou:2016jqi,Zhang:2017bzy,Monahan:2017hpu,Chen:2017lnm,Stewart:2017tvs,Izubuchi:2018srq,Alexandrou:2018pbm,Chen:2018xof,Chen:2018fwa,Liu:2018uuj,Bali:2018spj,Radyushkin:2018nbf,Lin:2018qky,Karpie:2018zaz}.

In this paper, we study the UV divergences of quasi-gluon operators to all orders in QCD perturbation theory. The UV structure of quasi-gluon operators could be much more complicated, as we will explain. We define general bare quasi-gluon operators as
\begin{align}\label{eq:ftgx}
\begin{split}
{\cal O}_{bg}^{\mu \nu \rho\sigma}(\xi) = F^{\mu \nu}(\xi)\,\Phi^{(a)}(\{\xi,0\})\, F^{{\rho\sigma}}(0)\,,
\end{split}
\end{align}
where $\Phi^{(a)}(\xi,0)={\cal P}e^{-ig_s\int_0^{1} \xi\cdot A^{(a)}(r \xi)\,dr}$ is a path ordered gauge link in adjoint representation. To be definite, we assume $\xi_\mu$ along $z$-direction and introduce a unit vector $n^\mu = (0,0,0, 1)$, defining $v\cdot n \equiv v_z$ for any vector $v_\mu$.  Due to the dimensional derivative operator in $F^{\mu\nu}$, superficial power counting tells us that the vertex between gluon field strength and gauge link could be linearly UV divergent. By using a cutoff regularization, one-loop calculation in Refs.~\cite{Wang:2017qyg,Wang:2017eel} indeed shows uncanceled linear divergences for this vertex, which would make the multiplicative renormalization of quasi-gluon operators almost impossible.  For the comparison, the corresponding vertex of quasi-quark operators can be at most logarithmic UV divergent.

Another complication comes from operator mixing via UV renormalization. In principle, the general quasi-gluon operator in Eq.~\eqref{eq:ftgx} could have 36 independent operators after taking into account the antisymmetry of gluon field strength, and all of them could be mixed under renormalization.  In literature, quasi-gluon PDFs are constructed from linear combinations of ${\cal O}_{bg}^{\mu \nu \rho\sigma}$ by contracting them with some ``tensors''. In Refs.~\cite{Ji:2013dva} and \cite{Wang:2017eel} and the first lattice simulation \cite{Fan:2018dxu}, different definitions of quasi-gluon PDFs were obtained by contracting ${\cal O}_{bg}^{\mu \nu \rho\sigma}$ by $n_\mu n_\rho  g_{\nu i}g_\sigma^{i}$, $n_\mu n_\rho  g_{\nu \sigma}$ and $g_{\mu 0}g_{\rho0}  g_{\nu \sigma}$, respectively. Renormalization of these quasi-gluon PDFs is non-trivial due to the mixing of these operators.

In this paper, we first perform an explicit one-loop calculation of quasi-gluon operators of an asymptotic gluon of momentum $p$, defined as $\langle g(p)|{\cal O}_{bg}^{\mu \nu \rho\sigma}(\xi) | g(p) \rangle$.  We use dimensional regularization (DR) to regularize both logarithmic and linear UV divergences, which respectively appear as poles around $d=4$ and $d=4-1/n$ at $n$-loop order. We find that linear UV divergences of one-loop correction to the gluon-gauge-link vertex are canceled under DR, which makes the multiplicative renormalizability of quasi-gluon operators a possibility. We then explore all possible UV divergent topologies of higher order diagrams.  Using gauge invariance, we find that all linear UV divergences from the gluon-gauge-link vertex are canceled to all-orders in perturbation theory.
Then, we find that all of the 36 independent quasi-gluon operators can be multipliticatively renormalized without mixing with any other operators. Combining with our proof for quasi-quark operators in Ref.~\cite{Ishikawa:2017faj}, our work presented in this paper for quasi-gluon operators completes the proof of multipliticative renormalization of the quasi-parton operators.
